\documentclass[12pt]{article}
\usepackage[utf8]{inputenc}
\usepackage[T1]{fontenc}
\usepackage{amsmath,amssymb}
\usepackage{geometry}
\geometry{margin=2.5cm}

\begin{document}

\noindent The following are true properties:

\begin{enumerate}
\item[1.] $\big(a,(b,c)\big) = \big((a,b),c\big)$
\item[2.] $(ca,cb) = c(a,b)$.
\item[3.] if $(a,b)=1$ and $(a,c)=1 \Rightarrow (a,bc)=1$.

\noindent generalization

\[
(a,b_i)=1,\ i=\overline{1,n} \ \Rightarrow \ (a,\,b_1b_2\cdots b_n)=1.
\]

\item[4.] $a\mid bc$, $(a,b)=1 \ \Rightarrow \ a\mid c$

\item[5.] $\left.\begin{aligned} a&\mid c\\ b&\mid c\\ (a,b)&=1\end{aligned}\right\} \Rightarrow ab\mid c$

\noindent generalization: \quad $m_i\mid c$, \ $(m_i,m_j)=1 \ \Rightarrow \ m_1m_2\cdots m_n\mid c$

\item[6.] $(a,b)=1 \Rightarrow$ [continuation unclear in the original: the visible fragment reads ``$c,ab$''.]\footnote{property 6 is incomplete/unclear in the original -- it possibly involves $(ab,c)$, but the exact statement is not certain.}
\end{enumerate}

\bigskip
\noindent (1) evident.

\noindent $(b,c)=d$.\footnote{the exact role of this notation in the proof is unclear in the original.}

\bigskip
\noindent (3). \ $1=(a,b)=(a,c)$

\bigskip
\noindent (4) \ $a\mid bc \Rightarrow (a,bc)=a$

\[
(a,c) \overset{?}{=} \big((a,bc),c\big) = \big(a,(bc,c)\big)
\]
\[
\Rightarrow (a,c)=a \ \Rightarrow \ a\mid c
\]
\noindent\footnote{the middle algebraic step is very compressed in the original and cannot be reconstructed with full certainty -- please check it.}

\bigskip
\noindent (5). \ $(ab,c)=$ \ [the page breaks off here.]

\bigskip
\noindent Let $p$ be an irreducible natural number $\Rightarrow p$ is prime

\noindent\underline{Proof}: By contradiction: \ $p$ not prime $\Rightarrow \exists\,a,b\in\mathbb{N}$ such that

\[
p\mid ab,\quad p\nmid a,\quad p\nmid b
\]
\[
\Rightarrow (p,a)=1 \text{ and } (p,b)=1 \overset{(3)}{\Longrightarrow} (p,ab)=1
\]
\[
\Rightarrow p=1 \quad \times.
\]

\end{document}
